% ============================================================================
% TUGAS AKHIR - JALUR MAGANG DUA SEMESTER (MADUSEM)
% Template LaTeX laporan Tugas Akhir jalur Magang Dua Semester
% Program Studi D3 Rekayasa Perangkat Lunak Aplikasi
% Fakultas Ilmu Terapan, Universitas Telkom
%
% Sumber: Template TA Madusem_5.docx (milik Prodi D3RPLA, tidak ikut dipush ke repo).
% Catatan: Jalur ini hanya memiliki 1 (satu) Dosen Pembimbing.
% ============================================================================

\documentclass[a4paper,12pt,oneside]{report}

% --- Paket standar ---
\usepackage[a4paper,margin=3cm]{geometry}
\usepackage[onehalfspacing]{setspace} % spasi 1.5 untuk body
\usepackage{array}
\usepackage{booktabs}
\usepackage{tabularx}
\usepackage{longtable}
\usepackage{multirow}
\usepackage{multicol}
\usepackage{graphicx}
\usepackage{float}
\usepackage{caption}
\usepackage{titlesec}
\usepackage{fancyhdr}
\usepackage{enumitem}
\usepackage{hyperref}
\usepackage{url}
\usepackage{pdflscape} % halaman landscape (halaman tambahan)

\hypersetup{colorlinks=true,linkcolor=black,urlcolor=blue,citecolor=black}
\captionsetup{font=small,labelfont=bf,labelsep=period}
\graphicspath{{assets/}}

% --- Format judul bab ---
\titleformat{\chapter}[display]{\normalfont\huge\bfseries\filcenter}{BAB \thechapter}{0pt}{}
\titleformat{\section}{\normalfont\large\bfseries}{\thesection}{1em}{}
\titleformat{\subsection}{\normalfont\normalsize\bfseries}{\thesubsection}{1em}{}
\titleformat{\subsubsection}{\normalfont\normalsize\bfseries}{\thesubsubsection}{1em}{}

% --- Penomoran: Gambar/Tabel/Lampiran per bab ---
\usepackage{chngcntr}
\counterwithin{figure}{chapter}
\counterwithin{table}{chapter}
\renewcommand{\thefigure}{\arabic{chapter}.\arabic{figure}}
\renewcommand{\thetable}{\arabic{chapter}.\arabic{table}}
% Lampiran pakai counter sendiri
\newcounter{lampiran}
\renewcommand{\thelampiran}{\arabic{lampiran}}

% --- Header/footer ---
\pagestyle{fancy}
\fancyhf{}
\fancyhead[L]{\small\leftmark}
\fancyfoot[C]{\thepage}
\renewcommand{\headrulewidth}{0.4pt}

\begin{document}

% =====================================================================
% HALAMAN COVER
% =====================================================================
\begin{titlepage}
    \centering
    \vspace*{1cm}
    {\bfseries\Large JUDUL TUGAS AKHIR DALAM BAHASA INDONESIA\\(maksimal 3 baris)\par}
    \vspace{1cm}
    {\bfseries\Large JUDUL TUGAS AKHIR DALAM BAHASA INGGRIS\\(maximum 3 lines)\par}
    \vspace{2cm}

    Dokumen ini ditujukan untuk memenuhi persyaratan\\
    Mata Kuliah Tugas Akhir\\
    \textbf{Jalur Magang Dua Semester}\par
    \vspace{2cm}

    Disusun oleh,\\
    NIM - NAMA\par
    \vfill

    \textbf{PROGRAM STUDI}\\
    \textbf{D3 REKAYASA PERANGKAT LUNAK APLIKASI}\\
    \textbf{FAKULTAS ILMU TERAPAN}\\
    \textbf{UNIVERSITAS TELKOM}\\
    \textbf{BANDUNG}\\
    202\ldots\par
\end{titlepage}

% =====================================================================
% LEMBAR PERSEMBAHAN
% =====================================================================
\chapter*{Lembar Persembahan}
\addcontentsline{toc}{chapter}{LEMBAR PERSEMBAHAN}

\begin{center}
    Untuk orang tuaku,\\
    orang paling berjasa di seluruh dunia.
\end{center}

\vspace{3cm}
\begin{flushleft}
    Nama \#1\\
    Nama \#2
\end{flushleft}

% =====================================================================
% LEMBAR PENGESAHAN
% =====================================================================
\chapter*{Lembar Pengesahan}
\addcontentsline{toc}{chapter}{LEMBAR PENGESAHAN}

\noindent\textit{Pastikan kalimat dan penggalan judul sama persis dengan cover}

\begin{center}
    {\bfseries JUDUL TUGAS AKHIR DALAM BAHASA INDONESIA}\\[6pt]
    {\bfseries JUDUL TUGAS AKHIR DALAM BAHASA INGGRIS}
\end{center}

\vspace{1cm}
\begin{table}[h]
    \centering
    \begin{tabular}{p{4cm}p{8cm}}
        Penulis                 & Nama Lengkap Mahasiswa \\
        NIM                     & 670\ldots\ldots \\
        \multicolumn{2}{c}{} \\
        Dosen Pembimbing        & Nama Dosen Pembimbing \\
        NIP                     & 078\ldots\ldots \\
    \end{tabular}
\end{table}

\vspace{2cm}
\begin{center}
    Tanggal Pengesahan: 31 Juli 2025 (tanggal pengumpulan revisi)
\end{center}

% =====================================================================
% PERNYATAAN
% =====================================================================
\chapter*{Pernyataan}
\addcontentsline{toc}{chapter}{PERNYATAAN}

Dengan ini saya menyatakan bahwa:
\begin{enumerate}[label=\arabic*.]
    \item Karya tulis ini adalah asli dan belum pernah diajukan untuk mendapatkan
    gelar akademik (Ahli Madya, Sarjana, Magister dan Doktor), baik di Fakultas
    Ilmu Terapan Universitas Telkom maupun di perguruan tinggi lainnya;
    \item karya tulis ini murni gagasan, rumusan, dan penelitian saya sendiri,
    tanpa bantuan pihak lain, kecuali arahan tim pembimbing atau tim promotor
    atau penguji;
    \item dalam karya tulis ini tidak terdapat cuplikan karya atau pendapat yang
    telah ditulis atau dipublikasikan orang lain, kecuali secara tertulis dengan
    jelas dicantumkan sebagai acuan dalam naskah dengan menyebutkan nama
    pengarang dan dicantumkan dalam daftar pustaka;
    \item saya mengijinkan karya tulis ini dipublikasikan oleh Fakultas Ilmu
    Terapan Universitas Telkom, dengan tetap mencantumkan saya sebagai penulis;
    dan
    \item Pernyataan ini saya buat dengan sesungguhnya dan apabila pada kemudian
    hari terdapat penyimpangan dan ketidakbenaran dalam pernyataan ini maka saya
    bersedia menerima sanksi akademik berupa pencabutan gelar yang telah
    diperoleh karena karya tulis ini, serta sanksi lainnya sesuai norma yang
    berlaku di Fakultas Ilmu Terapan Universitas Telkom.
\end{enumerate}

\vspace{2cm}
\begin{flushright}
    Bandung, 31 Juli 2025 (tanggal pengumpulan revisi)\\[6pt]
    Pembuat pernyataan,\\[3cm]
    (tulis nama jelas di sini tanpa gelar)
\end{flushright}

% =====================================================================
% KATA PENGANTAR
% =====================================================================
\chapter*{Kata Pengantar}
\addcontentsline{toc}{chapter}{KATA PENGANTAR}

Berisi kata pengantar; dan ditulis dalam bahasa Indonesia baku.

\vspace{3cm}
\begin{flushright}
    Bandung, \textless\textless Bulan\textgreater\ \textless\textless Tahun\textgreater\textgreater\\[6pt]
    Penulis
\end{flushright}

% =====================================================================
% ABSTRAK
% =====================================================================
\chapter*{Abstrak}
\addcontentsline{toc}{chapter}{ABSTRAK}

Abstrak adalah representasi dari isi dokumen yang singkat dan tepat. Tujuan
abstrak adalah memudahkan pembaca mendapatkan informasi mengenai Tugas Akhir
yang dibuat penulis tanpa harus membaca seluruh isi dokumen serta menghemat
waktu pembaca. Abstrak ditulis dalam 1 paragraf saja, dengan panjang tidak
lebih dari 250 (dua ratus lima puluh) kata dengan spacing 1. Abstrak tidak
mencantumkan landasan teori. Abstrak harus dibuat dengan ringkas dan mampu
menggambarkan alasan atau pentingnya Tugas Akhir ini, tujuan yang hendak
dicapai (atau fitur-fitur utama yang ada dalam produk Tugas Akhir),
langkah-langkah metode pengerjaannya, tools yang digunakan, dan kesimpulan.

\vspace{1cm}
\noindent\textbf{Kata Kunci:} (keywords abstrak 5-6 kata/frasa)

% =====================================================================
% ABSTRACT
% =====================================================================
\chapter*{Abstract}
\addcontentsline{toc}{chapter}{ABSTRACT}

\begin{spacing}{1}
\textit{Write your abstract in English. Using spacing 1 and italic.\\
DO NOT USE ENGLISH TRANSLATOR SOFTWARE WITHOUT FURTHER EDITING}

\vspace{1cm}
\noindent\textbf{Keywords:} (abstract keywords 5-6 words/phrases)
\end{spacing}

% =====================================================================
% DAFTAR ISI / GAMBAR / TABEL / LAMPIRAN
% =====================================================================
\tableofcontents
\listoffigures
\listoftables

% Daftar Lampiran (manual karena bukan environment baku)
\chapter*{Daftar Lampiran}
\addcontentsline{toc}{chapter}{DAFTAR LAMPIRAN}
\begin{enumerate}[label=Lampiran \arabic*.,leftmargin=*]
    \item \ldots
    \item \ldots
\end{enumerate}

% =====================================================================
% BAB I PENDAHULUAN
% =====================================================================
\chapter{Pendahuluan}

\section{Latar Belakang}
Sub bab ini berisi alasan mengapa judul/topik ini diambil. Penulisan sub bab
ini dimulai dari menjelaskan kondisi yang dialami saat ini, diikuti dengan
masalah yang dialami saat ini, dan diakhiri dengan solusi yang ditawarkan untuk
mengatasi masalah tersebut. Latar belakang harus menjelaskan urgensi dari
masalah yang dialami dan didukung oleh data-data ilmiah yang dapat
dipertanggungjawabkan kebenarannya. Contoh: pencatatan data yang ada di
perusahaan X masih manual (menggunakan buku) dan data yang ada tersebar di
masing-masing cabang, sehingga dengan adanya Sistem Informasi ini, data dapat
terpusat pada kantor induk dan juga dapat diakses oleh seluruh cabang~\cite{ref1}.

\section{Rumusan Masalah dan Solusi}
Sub bab ini berisi fokus permasalahan yang dapat dirumuskan dari latar belakang
yang telah dituliskan pada poin sebelumnya. Penulisan sub bab ini umumnya
menyebutkan masalah yang perlu diteliti dalam kalimat tanya, tetapi dapat pula
dibuat dalam bentuk kalimat pernyataan atau paragraf dengan rincian berbentuk
numbering.

Setiap masalah perlu diiringi dengan solusi sebagai proses penyelesaian masalah
yang telah dipaparkan sebelumnya. Rencana solusi yang ditawarkan untuk
menyelesaikan masalah yang dihadapi di tempat kerja/magang harus dapat
menunjukkan implementasi keilmuan dari program studi.

\section{Tujuan}
Sub bab ini berisi alasan pekerjaan ini dilakukan dan apa saja jenis pekerjaan
dan tanggung jawab yang dikerjakan oleh peserta kerja/magang. Penulisan sub bab
ini umumnya dalam bentuk kalimat pernyataan atau paragraf dengan rincian
berbentuk numbering.

Tujuan menyatakan seberapa jauh output (sasaran terukur) yang hendak dicapai
sebagai solusi untuk menjawab masalah-masalah yang telah dirumuskan dalam sub
bab Rumusan Masalah.

\section{Batasan Masalah}
Dalam perumusan masalah dapat dijelaskan definisi, asumsi, dan lingkup yang
menjadi batasan TA. Contoh:

Batasan masalah dalam pembuatan aplikasi ini adalah:
\begin{itemize}
    \item Aplikasi diimplementasikan pada smartphone Android minimal versi Lollipop.
    \item Dan seterusnya\ldots
\end{itemize}

\section{Penjadwalan Kerja}
Sub bab ini berisi hal-hal apa yang akan dilakukan untuk mencapai solusi yang
telah ditawarkan pada poin sebelumnya. Jadwal Pelaksanaan ini akan menjadi
pedoman baik bagi mahasiswa maupun para pembimbing akan target-target pencapaian
sesuai timeline yang disepakati.

Satuan waktu yang digunakan sebagai timeline dapat disesuaikan dengan kebutuhan
masing-masing pekerjaan yang dilakukan. Satuan waktu yang sebaiknya digunakan
adalah satuan jam, hari, atau minggu.

Contoh jadwal Pelaksanaan dalam satuan waktu minggu:

\begin{table}[H]
    \caption{Tabel Pelaksanaan Kerja}
    \label{tab:jadwal-kerja}
    \centering
    \footnotesize
    \begin{tabular}{|c|p{3.2cm}|*{12}{c|}}
        \hline
        \multirow{2}{*}{No} & \multirow{2}{*}{Deskripsi Kerja} & \multicolumn{12}{c|}{Bulan} \\ \cline{3-14}
         & & \multicolumn{3}{c|}{1} & \multicolumn{3}{c|}{2} & \multicolumn{3}{c|}{3} & \multicolumn{3}{c|}{4} \\ \hline
        1 & Diskusi      & \multicolumn{12}{c|}{} \\ \hline
        2 & Perancangan  & \multicolumn{12}{c|}{} \\ \hline
        3 & Penilaian    & \multicolumn{12}{c|}{} \\ \hline
        4 & Penelitian   & \multicolumn{12}{c|}{} \\ \hline
    \end{tabular}
\end{table}

\noindent\textbf{Petunjuk umum penulisan:}
\begin{itemize}
    \item Gunakan Bahasa Indonesia baku.
    \item Bila mengambil dari referensi, tulis ulang (paraphrase), jangan copy-paste.
    \item Judul Tabel dan Gambar seperti yang disertakan di bawah. Gunakan
    fasilitas Caption dan cross-reference pada Ms.\ Word.
    \item Semua tulisan pada gambar harus dapat terbaca dengan jelas. Gambar
    harus proporsional dan hi-res.
    \item Semua gambar dan tabel harus disebutkan pada narasi. Contoh, Tabel~\ref{tab:data-pengujian}
    menampilkan data pengujian, sementara Gambar~\ref{fig:logo-telkom}
    menampilkan logo universitas.
    \item Referensi menggunakan standar IEEE. Semua Daftar Pustaka harus ditulis
    dengan benar (terutama pada First Name -- Last Name).~\cite{ref1,ref2,ref3,ref4,ref5}.
    Gunakan fasilitas References pada Ms.\ Word.
    \item Gambar atau tabel yang lebih besar dari satu halaman, dituliskan pada
    lampiran.
\end{itemize}

\begin{table}[H]
    \caption{Data Pengujian}
    \label{tab:data-pengujian}
    \centering
    \begin{tabular}{|l|l|l|}
        \hline
        \textbf{Header 1} & \textbf{Header 2} & \textbf{Header 3} \\ \hline
        Isi 1 & Isi 2 & Isi 3 \\ \hline
              &       &       \\ \hline
              &       &       \\ \hline
              &       &       \\ \hline
    \end{tabular}
\end{table}

\begin{figure}[H]
    \centering
    \fbox{\parbox{0.5\textwidth}{\centering\vspace{3cm}(Logo Telkom University)\vspace{3cm}}}
    \caption{Logo Telkom University}
    \label{fig:logo-telkom}
\end{figure}

% =====================================================================
% BAB II PROFIL ORGANISASI
% =====================================================================
\chapter{Profil Organisasi}
Sub bab ini berisi profil lengkap perusahaan seperti nama perusahaan, alamat
perusahaan, sejarah perusahaan, logo perusahaan dan hal-hal lain yang perlu
dijelaskan mengenai brand perusahaan. Apabila ada gambar yang perlu dicantumkan
wajib diiringi dengan penjelasan gambar tersebut. Termasuk juga struktur
organisasi institusi dicantumkan pada sub bab ini.

\section{Deskripsi Organisasi}
Bagian ini berisi penjelasan menggambarkan organisasi lokasi magang secara umum,
meliputi nama perusahaan, merk/brand startup beserta filosofi logo, bidang,
produk yang dimiliki, tanggal pendirian atau tanggal inisiasi pembentukan
startup, ijin usaha (jika ada), bentuk perusahaan saat ini, dan
kepemilikan/tim pendiri. Serta deskripsi usaha yang dilaksanakan oleh organisasi
secara singkat.

\section{Struktur Organisasi dan Tata Kelola}
Sub bab ini berisi penjelasan spesifik mengenai unit kerja, divisi kerja, atau
bagian kerja. Beberapa hal yang perlu disampaikan dalam sub bab ini diantaranya
adalah:
\begin{enumerate}
    \item Deskripsi mengenai divisi tempat kerja/magang;
    \item Daftar pekerjaan secara umum yang dilakukan oleh divisi tempat kerja/magang;
    \item Keterkaitan pekerjaan tersebut dengan divisi lain;
    \item Peran kerja divisi terhadap proses bisnis perusahaan secara utuh.
\end{enumerate}
Penulisan sub bab ini umumnya dalam bentuk narasi atau dengan pemodelan tertentu
(misalnya flowchart).

\section{Deskripsi Pekerjaan}
Sub bab ini berisi alasan pekerjaan ini dilakukan dan apa saja jenis pekerjaan
dan tanggung jawab yang dikerjakan oleh peserta kerja/magang. Berilah keterangan
seberapa jauh ruang lingkup pekerjaan yang dikerjakan oleh peserta kerja/magang.

Juga ditambahkan penjelasan alur kerja yang dilakukan oleh peserta kerja/magang
dimulai dari persiapan hingga pekerjaan dianggap selesai dan/atau dilimpahkan
pada unit lainya. Penulisan sub bab ini umumnya dalam bentuk pemodelan tertentu
(misalnya flowchart) yang dijelaskan ke dalam bentuk narasi paragraf.

% =====================================================================
% BAB III ANALISIS PEKERJAAN
% =====================================================================
\chapter{Analisis Pekerjaan}
Bagian ini berisi aktivitas yang dilakukan selama mengikuti program magang 2
semester, dan capaian yang telah diperoleh melalui aktivitas tersebut.

\section{Tinjauan Pustaka}
Tinjauan pustaka berisi teori-teori yang akan digunakan dalam pengerjaan tugas
akhir. Jadi bukan tentang sejarah sebuah framework atau tools tertentu. Sebagai
contoh, ketika membahas tentang Android, fokus bahasan bukan tentang Android
dibuat oleh siapa, diakuisisi sama siapa dan kapan kejadiannya. Tapi bahaslah
tentang komponen aplikasi Android seperti activity, services, broadcast receiver
dan content provider yang mungkin akan dipakai di aplikasi kalian nantinya.

\section{Analisis Sistem}
Sebagai sebuah solusi yang dapat terukur, penulis harus menentukan
ukuran/kriteria kinerja dan menjadi acuan kualitas sistem yang dikembangkan.

\subsection{Gambar Sistem Saat Ini}
Sub bab ini berisi penjelasan mengenai sistem atau alur kerja yang ada saat ini
dialami oleh peserta kerja/magang. Penulisan sub bab ini umumnya dalam bentuk
paragraf narasi.

\subsection{Pengembangan Sistem}
Sub bab ini berisi penjelasan mengenai pengembangan/perubahan sistem atau alur
kerja yang diajukan oleh peserta kerja/magang. Penulisan sub bab ini umumnya
dalam bentuk paragraf narasi atau diagram. Sebagai contoh:

\begin{itemize}
    \item \textbf{Pada bidang Computer Technical Support, Network Administrator,
    System Administrator.} Dapat berupa diagram sistem (baik sistem maupun
    jaringan) serta file konfigurasi.
    \item \textbf{Pada bidang Programmer, Web Programmer, Mobile Programmer.}
    Dapat berupa perangkat lunak yang dibangun (mockup antar muka), bisa desain
    sebelum implementasi maupun hasil setelah implementasi. Juga bisa dimasukkan
    diagram rancangan perangkat lunak.
    \item \textbf{Pada bidang Database Programmer.} Dapat berupa query, tabel,
    database, subprogram dan pengimplementasiannya. File source .sql dapat
    dilampirkan juga beserta kumpulan code query yang dibangun.
    \item \textbf{Pada bidang Webmaster dan Multimedia Programmer.} Dapat berupa
    alur penggunaan konten/multimedia yang dibangun (dapat dinyatakan dengan
    flowchart ataupun diagram bisnis proses).
    \item \textbf{Pada bidang Embedded Programmer.} Dapat berupa diagram
    rancangan maupun alur kerja hardware dan juga diagram desain hardware yang
    dibangun.
    \item \textbf{Pada Bidang Quality Assurance.} Dapat berupa dokumen pengujian
    dan skema pengujian yang anda bangun. Serta langkah-langkah penggunaan tools
    pengujian yang digunakan dalam pengujian.
\end{itemize}

\section{Kebutuhan Perangkat Kerja}
Sub bab ini berisi penjelasan perangkat kerja yang dipakai oleh peserta
kerja/magang dimulai dari persiapan hingga pekerjaan dianggap selesai dan/atau
dilimpahkan pada unit lainya. Penulisan sub bab ini umumnya dalam bentuk
numbering atau tabel yang menjelaskan perangkat keras dan perangkat lunak apa
saja yang digunakan dalam menyelesaikan pekerjaan. Jelaskan fungsi dari perangkat
keras dan perangkat lunak yang digunakan dan lampirkan gambar atau screenshot
dari perangkat tersebut.

\subsection{Kebutuhan Perangkat Pengembangan Sistem}
Membahas kebutuhan perangkat keras dan perangkat lunak yang digunakan dalam
pengembangan sistem.

\subsection{Kebutuhan Perangkat Implementasi Sistem}
Membahas kebutuhan perangkat keras dan perangkat lunak yang digunakan dalam
implementasi sistem yang telah dikembangkan di lingkungan operasional.

\subsection{Sub Bab Tambahan (Bila Diperlukan)}
Mohon disesuaikan.

% =====================================================================
% BAB IV HASIL DAN PEMBAHASAN
% =====================================================================
\chapter{Hasil dan Pembahasan}
Bagian ini berisi deskripsi atau narasi yang menjelaskan hasil pengembangan
sistem yang menjawab permasalahan pada organisasi lokasi magang.

\section{Hasil Akhir (Luaran)}
Sub bab ini berisi penjelasan mengenai hasil akhir (luaran) dari kerja/magang
yang telah dilaksanakan secara umum dan proses selanjutnya setelah luaran selesai
secara umum. Penulisan sub bab ini pada umumnya berbentuk paragraf narasi.

Sub bab ini berisi mengenai detail hasil akhir dari kerja/magang. Penulisan dan
isi dari bagian ini akan berbeda, sesuai dengan bidang kerja selama
kerja/magang. Berikut contohnya.

\begin{itemize}
    \item \textbf{Pada bidang Computer Technical Support, Network Administrator,
    System Administrator.} Luarannya dapat berupa diagram sistem (baik sistem
    maupun jaringan) serta semua file konfigurasi dari sistem yang sudah dibangun
    dan diatur (contoh file konfigurasi adalah file \texttt{httpd.conf} pada
    apache web server). Bisa dicantumkan juga hasil pengujian sistem dan
    implementasi. Foto dokumentasi sistem yang anda bangun dan jaringannya juga
    dilampirkan.
    \item \textbf{Pada bidang Programmer, Web Programmer, Mobile Programmer.}
    Luarannya dapat berupa perangkat lunak yang dibangun (mockup antar muka).
    Selain itu, apabila diperlukan source code dapat dilampirkan di laporan
    akhir. Jika source code terlalu besar maka dapat dilampirkan dalam sebuah
    CD/DVD. Printscreen aplikasi/web juga dilampirkan. Bisa dicantumkan juga
    hasil pengujian sistem dan implementasi perangkat lunak.
    \item \textbf{Pada bidang Database Programmer.} Luarannya dapat berupa query,
    tabel, database, subprogram dan pengimplementasiannya. File source .sql
    dapat dilampirkan juga beserta kumpulan code query yang dibangun. Bisa
    dicantumkan juga hasil pengujian database.
    \item \textbf{Pada bidang Webmaster dan Multimedia Programmer.} Luarannya
    dapat berupa file konten/multimedia juga dapat dilampirkan berupa
    printscreennya konten yang dibangun. File multimedia dapat dilampirkan dalam
    sebuah CD/DVD. Bisa dicantumkan juga hasil pengujian aplikasi multimedia.
    \item \textbf{Pada bidang Embedded Programmer.} Luarannya dapat berupa source
    code dan hardware. File source code dapat dilampirkan di laporan ataupun
    CD/DVD jika dirasa terlalu besar. Hardware yang dibangun juga dapat difoto.
    Bisa dicantumkan juga hasil pengujian sistem dan implementasinya.
    \item \textbf{Pada bidang Quality Assurance.} Luarannya dapat berupa dokumen
    pengujian dan skema pengujian yang dibangun. Serta langkah-langkah penggunaan
    tools pengujian yang digunakan dalam pengujian. Beberapa printscreen pengujian
    yang dilakukan juga dilampirkan. Juga bisa dicantumkan rencana perbaikan dari
    hasil pengujian.
\end{itemize}

\section{Pengujian Luaran}
Membahas pengujian dari sistem yang telah dikembangkan atau luaran lain yang
sesuai dengan permasalahan yang diselesaikan. Teknik pengujian dapat
disesuaikan dengan karakteristik luaran yang telah dihasilkan.

\subsection{Sub Bab Lain yang Relevan (Opsional)}
Mohon disesuaikan.

% =====================================================================
% BAB V PENUTUP
% =====================================================================
\chapter{Penutup}
Bagian ini berisi kesimpulan dari solusi yang telah dikembangkan untuk
organisasi lokasi magang serta saran pengembangan kedepannya.

\section{Kesimpulan}
Sub bab ini berisi ringkasan pekerjaan mulai dari hal-hal yang memotivasi
pekerjaan atau penelitian sampai ke hasil pekerjaan. Harap diperhatikan bahwa
kesimpulan harus menjawab tujuan dan bersifat ilmiah. Kesimpulan adalah
pembuktian bahwa tujuan dari pekerjaan atau penelitian yang telah tercapai.
Jangan pernah menuliskan apa yang tidak pernah Anda analisis atau kerjakan
sebelumnya.

\section{Saran}
Sub bab ini berisi hal-hal baru seperti alat/aplikasi/hasil survey yang
diperoleh penulis yang dapat memperbaiki atau mengembangkan esensi dari
pembahasan. Sertakan alasannya. Sebagai contoh, disarankan untuk menggunakan
database Oracle karena mampu menangani data dalam jumlah yang lebih besar
dibandingkan MySQL; dari analisis survey disarankan untuk menurunkan harga
makanan sehingga lebih terjangkau bagi mahasiswa.

% =====================================================================
% DAFTAR PUSTAKA
% =====================================================================
\renewcommand{\bibname}{Daftar Pustaka}
\begin{thebibliography}{9}
\bibitem{ref1}
I. Tin, ``Liputan 6 Tekno,'' 21 November 2016. [Online]. Available:
\url{http://tekno.liputan6.com/read/2591339/tak-disangka-pengguna-bitcoin-di-indonesia-capai-200-ribu}.

\bibitem{ref2}
\textit{(Referensi kedua -- tulis sesuai format IEEE)}

\bibitem{ref3}
\textit{(Referensi ketiga -- tulis sesuai format IEEE)}

\bibitem{ref4}
\textit{(Referensi keempat -- tulis sesuai format IEEE)}

\bibitem{ref5}
\textit{(Referensi kelima -- tulis sesuai format IEEE)}
\end{thebibliography}

% =====================================================================
% LAMPIRAN
% =====================================================================
\appendix
\chapter{Lampiran}
\stepcounter{lampiran}
Bagian ini berisi lampiran-lampiran pendukung laporan tugas akhir.

\section*{Lampiran 1}
\addcontentsline{toc}{section}{Lampiran 1}
\textit{(Isi lampiran 1)}

\section*{Lampiran 2}
\addcontentsline{toc}{section}{Lampiran 2}
\textit{(Isi lampiran 2)}

% =====================================================================
% HALAMAN TAMBAHAN (contoh landscape)
% =====================================================================
\newpage
\begin{landscape}
\chapter*{Halaman Tambahan}
\addcontentsline{toc}{chapter}{HALAMAN TAMBAHAN}
\noindent\textit{Note: ini adalah contoh untuk halaman Landscape. Hapus jika tidak diperlukan.}
\end{landscape}

\end{document}
